% 李群（综述）
% license CCBYSA3
% type Wiki

本文根据 CC-BY-SA 协议转载翻译自维基百科\href{https://en.wikipedia.org/wiki/Lie_group}{相关文章}

在数学中，李群（Lie group，发音为 /liː/，读作“LEE”）是既是群又是可微流形的结构，其特点是群的乘法运算和取逆运算都是可微的。

流形是局部上类似于欧几里得空间的空间；而群则定义了一种抽象的二元运算及其必须具备的附加性质，使其能被看作一种“抽象变换”。例如，群可以通过乘法与逆元运算（对应于除法），或等价地，通过加法与减法来体现。将这两种概念结合，就得到了连续群，其中点的乘法和取逆运算是连续的。如果乘法和取逆进一步是光滑（可微）的，那么这样的连续群就是李群。

李群为连续对称性的概念提供了一种自然模型，其中著名的例子就是圆群。旋转一个圆就是连续对称性的实例。对于圆的任意一次旋转，都存在同样的对称性；而这些旋转的组合（拼接）构成了圆群，这是李群的典型例子。李群在现代数学和物理的许多领域都有广泛的应用。

李群最初是在研究矩阵子群
$G \subseteq \mathrm{GL}_n(\mathbb{R})$ 或 $\mathrm{GL}_n(\mathbb{C})$ 时被发现的，这里的 $\mathrm{GL}_n(\mathbb{R})$ 与 $\mathrm{GL}_n(\mathbb{C})$ 分别表示实数域 $\mathbb{R}$ 或复数域 $\mathbb{C}$ 上的 $n \times n$ 可逆矩阵群。如今，这些群被称为**经典群**，因为李群的概念已经远远超出了这些最初的范围。

李群的名称源自挪威数学家索弗斯·李（Sophus Lie，1842–1899），他奠定了连续变换群理论的基础。李引入李群的最初动机，是为了刻画微分方程的连续对称性，这与伽罗瓦理论中利用有限群来刻画代数方程的离散对称性有着类似的思想。
\subsection{历史}
索弗斯·李认为1873–1874 年冬天是他连续群理论的诞生时期。[2] 然而，托马斯·霍金斯则认为，应当追溯到“李在 1869 年秋天至 1873 年秋天这四年间的惊人研究活动”，这才真正促成了该理论的建立。[2]李的一些早期思想是在与费利克斯·克莱因的密切合作中发展起来的。从 1869 年 10 月到 1872 年，李几乎每天都与克莱因见面：从 1869 年 10 月底到 1870 年 2 月底在柏林，其后两年则在巴黎、哥廷根和埃尔兰根。[3]李本人声称，到 1884 年时，他的主要成果都已经获得。然而，在 1870 年代，他的所有论文（除了最初的一篇简短笔记）都发表在挪威的学术期刊上，这阻碍了这些成果在欧洲其他地区的传播与认可。[4]1884 年，一位年轻的德国数学家弗里德里希·恩格尔来到李身边，与他一起撰写一部系统性著作，以全面阐述连续群理论。这一合作的成果就是三卷本的《变换群论》，分别于 1888 年、1890 年和 1893 年出版。“李群”这一术语最早于1893 年 出现，在李的学生阿图尔·特雷斯的博士论文中首次使用。[5]

李的思想并非孤立于数学整体之外。事实上，他对微分方程几何的兴趣最初是受到卡尔·古斯塔夫·雅可比工作的启发——雅可比曾研究一阶偏微分方程理论以及经典力学中的方程。雅可比的许多成果在他去世后的 1860 年代才得以发表，这在法国和德国引起了极大的兴趣。[6]李的执念（是建立一个微分方程对称性的理论，使其能完成与伽罗瓦在代数方程中所做的相同使命：即利用群论对它们进行分类。李和其他数学家证明，许多重要的特殊函数与正交多项式的方程，往往源于群论对称性。在李的早期工作中，他的目标是建立一个连续群的理论，以补充费利克斯·克莱因与昂利·庞加莱在模形式理论中所发展的离散群理论。李最初心中所想的应用领域，是微分方程理论。借鉴伽罗瓦理论及其在多项式方程中的作用，他的核心构想是：通过研究对称性，建立一个能够统一常微分方程整体领域的理论。然而，这一愿望并未完全实现。虽然对称性方法在常微分方程（ODE）中仍然是一个活跃的研究方向，但它并未成为该领域的主导工具。后来出现了微分伽罗瓦理论，但它是由其他数学家（如皮卡尔 Picard 和 维西奥 Vessiot）发展起来的，主要提供了一种求积理论，即通过不定积分来表达解的框架。

进一步推动人们思考连续群的动力，来自伯恩哈德·黎曼关于几何基础的思想，以及这些思想在克莱因（Klein）手中的进一步发展。由此，李在创建他的新理论时，结合了 19 世纪数学中的三大主题：

\begin{itemize}
\item 对称性的思想 —— 伽罗瓦通过群这一代数概念所展现的典范；
\item 几何理论与力学微分方程的显式解 —— 由泊松与雅可比所发展；
\item 几何学的新理解 —— 在普吕克、莫比乌斯、格拉斯曼等人的研究中逐步形成，并在黎曼的革命性愿景中达到顶峰。
\end{itemize}
虽然今天索弗斯·李被公认为连续群理论的奠基人，但在其结构理论的发展上，一个具有决定性意义的飞跃来自威廉·基林。他在1888 年发表了一系列题为《连续有限变换群的合成》 的论文中的第一篇，这一工作对后续数学的发展产生了深远影响。[7]

基林的成果后来被埃利·卡当加以精炼与推广，最终推动了半单李代数的分类、卡当的对称空间理论，以及赫尔曼·外尔利用最高权描述紧李群与半单李群的表示理论的建立。

1900年，大卫·希尔伯特在巴黎召开的国际数学家大会上提出了著名的希尔伯特第五星问题，向李群理论的研究者发出了挑战。

外尔推动了李群理论早期发展的成熟。他不仅完成了半单李群不可约表示的分类，并且将群论与量子力学建立了紧密联系；更重要的是，他通过清晰阐明李的无穷小群（即李代数）与李群本身之间的区别，使李的理论基础更加稳固，并开启了对李群拓扑结构的研究。[8]李群理论随后在现代数学语言的框架下被系统地重新表述，这一工作体现在克洛德·谢瓦莱所撰写的专著中。
\subsection{概述}
\begin{figure}[ht]
\centering
\includegraphics[width=6cm]{./figures/38f113f77dfdff69.png}
\caption{所有模为1的复数集合（对应于复平面上以原点为中心、半径为1的圆上的点），在复数乘法下构成一个李群，这就是圆群。} \label{fig_LQun_1}
\end{figure}
李群是光滑的可微流形，因此可以利用微积分进行研究，这与更一般的拓扑群情况形成对比。李群理论中的一个关键思想，是用群的局部或线性化版本来取代整体对象——李本人称之为其“无穷小群”，后来被称为李代数。

李群在现代几何中具有极其重要的作用，并且体现在多个层面。费利克斯·克莱因在他的埃尔兰根纲领中指出，可以通过指定一个保持某些几何性质不变的合适变换群来考虑不同的“几何”。因此：欧几里得几何对应于欧几里得空间 $\mathbb{R}^3$ 上保持距离不变的变换群 $E(3)$；共形几何对应于将该群扩展到共形群；射影几何则关注在射影群下保持不变的性质。这一思想后来引出了G-结构的概念，其中$G$是流形的“局部”对称性李群。

李群（及其相关的李代数）在现代物理学中发挥着重要作用，通常作为物理系统的对称性出现。在这里，李群（或其李代数）的表示具有特别重要的意义。表示论在粒子物理学中被广泛应用。其中特别重要的群包括：旋转群$\mathrm{SO}(3)$（或其双覆盖群 $\mathrm{SU}(2)$），特殊酉群$\mathrm{SU}(3)$，庞加莱群。

在“整体”层面上，每当一个李群作用于某个几何对象（例如黎曼流形或辛流形）时，这种作用都会带来某种刚性，并产生丰富的代数结构。通过李群在流形上的作用来表达的连续对称性的存在，会对该流形的几何结构施加强约束，并有助于在其上进行分析。李群的线性作用尤为重要，它们是表示论的核心研究对象。

在 1940–1950 年代，埃利斯·科尔钦、阿尔芒·博雷尔以及克洛德·谢瓦莱认识到，许多关于李群的基础性结果可以完全用代数方法来建立，从而产生了代数群理论，该理论允许在任意域上定义代数群。这一洞见为纯代数开辟了新的可能性：它不仅提供了一种对大多数有限单群的统一构造，也在代数几何中发挥了深远作用。自守形式理论作为现代数论的重要分支，与李群的类似物（定义在阿德尔环上的群）密切相关；而p进李群则因其与数论中伽罗瓦表示的联系，而扮演着关键角色。
\subsection{定义与例子}
一个实李群是既是群又是有限维实光滑流形的对象，其中群运算（乘法与取逆）都是光滑映射。群乘法
$$
\mu : G \times G \to G, \quad \mu(x,y) = xy~
$$
的光滑性意味着$\mu$是从乘积流形$G \times G$到$G$的光滑映射。

这两个要求（乘法和取逆的光滑性）可以合并为一个条件：映射
$$
(x,y) \mapsto x^{-1}y~
$$
是从乘积流形$G \times G$到$G$的光滑映射。
\subsubsection{首批例子}

所有 $2 \times 2$ 实可逆矩阵在矩阵乘法下构成一个群，称为**2 阶一般线性群**，记作

$$
\operatorname{GL}(2,\mathbb{R}) \quad \text{或} \quad \operatorname{GL}_{2}(\mathbb{R}) :
$$

$$
\operatorname{GL}(2,\mathbb{R})=\left\{A=\begin{pmatrix}a & b \\ c & d\end{pmatrix} : \det A = ad - bc \neq 0 \right\}.
$$

这是一个 **四维非紧实实李群**；它是 $\mathbb{R}^4$ 的一个开子集。该群是不连通的，具有两个连通分支，分别对应于行列式取正值和负值的情况。

**旋转矩阵**构成 $\operatorname{GL}(2,\mathbb{R})$ 的一个子群，记作

$$
\operatorname{SO}(2,\mathbb{R}).
$$

它本身就是一个李群：具体来说，是一个**一维紧致连通李群**，并且与圆光滑同胚。若用旋转角 $\varphi$ 作为参数，则该群可参数化为：

$$
\operatorname{SO}(2,\mathbb{R})=\left\{\begin{pmatrix}\cos \varphi & -\sin \varphi \\ \sin \varphi & \cos \varphi \end{pmatrix} : \varphi \in \mathbb{R} / 2\pi\mathbb{Z}\right\}.
$$

角度的加法对应于 $\operatorname{SO}(2,\mathbb{R})$ 中元素的乘法，而取相反角度对应于取逆运算。因此，群的乘法与取逆运算都是可微映射。

**一维仿射群**是一个二维矩阵李群，由所有满足对角元分别为正数和 1 的 $2 \times 2$ 实上三角矩阵组成：

$$
A=\begin{pmatrix} a & b \\ 0 & 1 \end{pmatrix}, \quad a > 0, \, b \in \mathbb{R}.
$$
